% Inactive reconstruction. Not an author input or a release certificate.
\subsection{Curve maps with critical fibers}
\label{subsec:curve-map-critical-fibers}

Let $p:X\to C$ be a surjective holomorphic map with connected
fibers, where $X$ is a connected compact complex manifold of
dimension $n\geq1$ and $C$ is a connected compact Riemann surface.
The map need not be a submersion. In the first five results,
$\beta$ is any smooth real $(1,1)$-form on $C$.

\begin{proposition}\label{curve-map:potential-descent}
Pullback gives a bijection
\[
 \PSH(C,\beta)\longrightarrow\PSH(X,p^*\beta).
\]
It preserves the supremum and the ordinary singularity order.
\end{proposition}
\begin{proof}
The critical set is a proper compact analytic subset of $X$.
It has finitely many irreducible components, and $p$ is constant
on each: its differential vanishes on the smooth locus of that
component, and continuity extends the constant to the component.
Thus the set $D\subset C$ of critical values is finite.

Let $U\in\PSH(X,p^*\beta)$. On a base coordinate disk choose a
smooth $a$ with $\ddc a=\beta$. Over $C\setminus D$, the function
$U+a\circ p$ restricts to a plurisubharmonic function, or the
constant $-\infty$, on each compact connected fiber. The maximum
principle makes it constant there. Local holomorphic sections
then give a function $v$ with $U=v\circ p$ and $a+v$ subharmonic.
This is the local argument of \cref{prop:psh-descent-submersion};
it uses compactness of the fibers, not compactness of the
punctured total space.

The function $v$ is bounded above by $\sup_XU$. The removable
singularity theorem extends $a+v$ subharmonically across the
finitely many punctures. The extensions glue, giving
$v\in\PSH(C,\beta)$. Its pullback is $p^*\beta$-plurisubharmonic
by \cref{lem:smooth-disk-criterion}. It agrees with $U$ off the
proper analytic set $p^{-1}(D)$, hence almost everywhere and
everywhere by \cref{lem:qpsh-almost-everywhere-order}.

Conversely pullback preserves the curvature inequality by the
disk criterion. It is not identically $-\infty$, because $p$ is
surjective. Surjectivity also proves uniqueness and preserves
the supremum and every inequality $v\leq w+c$.
\end{proof}

\begin{lemma}\label{curve-map:lelong-pullback}
Let $f:(\mathbb C^m,0)\to(\mathbb C,0)$ be a nonconstant
holomorphic germ, with $m\geq1$, and let $e=\operatorname{ord}_0f$.
For every quasi-subharmonic germ $v$ at zero,
\[
 \nu(v\circ f,0)=e\nu(v,0).
\]
\end{lemma}
\begin{proof}
A smooth correction of $v$ pulls back to a smooth function, so
we may assume that $v$ is subharmonic. The Taylor expansion of
$f$ gives $|f(z)|\leq M|z|^e$ on a small ball. Choose a unit vector
$\xi$ on which the leading homogeneous term is nonzero. Then
\[
 f(\zeta\xi)=a\zeta^e+O(\zeta^{e+1}),\qquad a\ne0.
\]
For all small $r$, the remainder has modulus at most
$|a|r^e/4$ on $|\zeta|=r$. If $|t|<|a|r^e/4$, Rouch\'e's theorem
shows that $f(\zeta\xi)-t$ has $e$ zeros in $|\zeta|<r$.
Consequently, with $c=|a|/4$,
\[
 \{|t|<cr^e\}\subset f(B(0,r))
                  \subset\{|t|<Mr^e\}.
\]
The suprema of $v$ on these three sets are ordered. Dividing
by $\log r^2<0$, the two disk terms both tend to $e\nu(v,0)$
by the definition in \cref{sec:analytic-types}. The middle term
therefore has the same limit, proving the formula.
\end{proof}

\begin{corollary}\label{curve-map:ideal-types}
For $v,w\in\QPSH(C)$, the inclusions
$\mathcal I(tv)\subseteq\mathcal I(tw)$ for every $t>0$ hold
if and only if
\[
 \mathcal I(t(v\circ p))\subseteq\mathcal I(t(w\circ p))
 \quad\text{for every }t>0.
\]
In particular pullback preserves and reflects all-weight
multiplier-ideal equivalence.
\end{corollary}
\begin{proof}
On a smooth curve, local modifications are isomorphisms.
By \cref{thm:local-i-inputs}, the base inclusions are therefore
equivalent to $\nu(v,s)\geq\nu(w,s)$ for every $s\in C$.
Let $\pi$ be any local modification of a neighborhood in $X$
with smooth source, and let $y$ be a point of its source.
In a coordinate at $s=p(\pi(y))$, the germ $p\circ\pi-s$ is
nonconstant. Indeed, a modification is an isomorphism on a
dense open set, and a nonconstant holomorphic map cannot be
constant on an open set. Its order $e_y$ is positive and finite.
The preceding lemma gives
\[
 \nu(v\circ p\circ\pi,y)=e_y\nu(v,s),\qquad
 \nu(w\circ p\circ\pi,y)=e_y\nu(w,s).
\]
The base inequalities thus imply all inequalities upstairs,
and the same valuative theorem gives the ideal inclusions.
Conversely the upstairs inclusions give the inequalities for
the identity modification at any $x\in p^{-1}(s)$. Dividing by
$\operatorname{ord}_x(p-s)>0$ recovers the base inequalities.
Apply this argument in both directions for equality.
\end{proof}

\begin{lemma}\label{curve-map:lone-pullback}
If $v_j,v\in\PSH(C,\beta)$ and $v_j\to v$ in $L^1(C)$, then
$v_j\circ p\to v\circ p$ in $L^1(X)$ for every smooth positive
volume form on $X$.
\end{lemma}
\begin{proof}
Put $M_j=\sup_Cv_j$ and $M=\sup_Cv$. Hartogs' lemma applied at
moving maximum points gives $\limsup M_j\leq M$. A subsequence
with limiting suprema below $M$ would give $v$ below that limit
almost everywhere, hence everywhere, a contradiction. Thus
$M_j\to M$.

By surjectivity, $(v_j-M_j)\circ p$ has supremum zero. Normalized
compactness, \cref{thm:normalized-compactness}, makes these
pullbacks precompact in $L^1(X)$. On every relatively compact
product chart over $C\setminus D$, where $D$ is the finite set
of critical values, Fubini's theorem and bounded smooth volume
densities give local $L^1$ convergence to $(v-M)\circ p$.
Every subsequential limit therefore equals this pullback off
$p^{-1}(D)$, and hence almost everywhere, because that analytic
set has smooth volume zero. All subsequential limits agree,
so the entire normalized sequence converges. Adding $M_j\to M$
proves the assertion.
\end{proof}

\begin{lemma}\label{curve-map:analytic-types}
A function $v\in\QPSH(C)$ has analytic singularities if and
only if $v\circ p$ has analytic singularities.
\end{lemma}
\begin{proof}
Composing a local analytic presentation with $p$ proves the
forward implication. The pulled-back generators are not all
identically zero, since a nonconstant map to a curve is open.

Conversely fix $s\in C$ and $x\in p^{-1}(s)$. In coordinates,
choose a complex line through $x$ on which $p-s$ has its finite
order $e$ and is nonconstant, as in the Lelong lemma. On that
line write
\[
 (p-s)(\zeta)=h(\zeta)^e,
\]
where $h$ is a local biholomorphism fixing zero; take an analytic
$e$th root of the nonvanishing factor of the germ. The restriction
of $v\circ p$ to the line is not identically $-\infty$, since
$v$ is finite almost everywhere and this map is open. By
\cref{prop:analytic-type-calculus}, it has analytic singularities.
In one variable its presentation is
$r\log|\zeta|^2+O(1)$ with $r\in\mathbb Q_{\geq0}$, obtained by
taking the least order of its finitely many generators.
Changing to $w=h(\zeta)$ gives
\[
 v(s+w^e)=r\log|w|^2+O(1).
\]
Every sufficiently small $t$ has an $e$th root $w$, so
$v(s+t)=(r/e)\log|t|^2+O(1)$. This is an analytic presentation
on the base, including a bounded one when $r=0$.
\end{proof}

\subsection{Gauduchon traces of critical fibers}

Fix a Gauduchon form $\gamma$ on $X$ and a smooth positive
area-one form $\eta$ on $C$. For a finite positive measure
$\mu$ of mass $M$, let $q_\mu$ be the solution with mean zero to
\[
 \ddc q_\mu=\mu-M\eta.
\]
The curve Green kernel of \cref{sec:curve-trace-realization}
gives this quasi-subharmonic solution. Define the positive current
\[
 p^*\mu=M p^*\eta+\ddc(q_\mu\circ p),\qquad
 L_\gamma(\mu)=p^*\mu\wedge\gamma^{n-1}.
\]
Locally this is the pullback of a subharmonic potential of $\mu$;
changing that potential adds a harmonic function and does not
change the current. Extend both operators to signed measures by
linearity. Write $F_s=p^{-1}(s)$ for the fiber divisor with its
scheme multiplicities.

\begin{lemma}\label{curve-map:gauduchon-trace}
There is a positive constant $\kappa_\gamma$ such that, for every
finite signed measure $\mu$ and every Borel set $B\subset C$,
\[
 \begin{gathered}
 p_*L_\gamma(\mu)=\kappa_\gamma\mu,\qquad
 \|L_\gamma(\mu)\|_{\mathrm{TV}}
               =\kappa_\gamma\|\mu\|_{\mathrm{TV}},\\
 L_\gamma(\mathbf1_B\mu)
            =\mathbf1_{p^{-1}(B)}L_\gamma(\mu),\qquad
 \int_{F_s}\gamma^{n-1}=\kappa_\gamma\quad(s\in C).
 \end{gathered}
\]
The last integral includes the multiplicities of the fiber.
\end{lemma}
\begin{proof}
The positive $(0,0)$-current $S=p_*(\gamma^{n-1})$ on $C$ satisfies
$\ddc S=p_*\ddc(\gamma^{n-1})=0$. A harmonic distribution is smooth,
so $S$ is a constant $\kappa_\gamma$. On the complement of the
critical values, smooth fiber integration identifies it with
the positive integral of $\gamma^{n-1}$ on a regular fiber.
Thus $\kappa_\gamma>0$.

First let $\mu$ be positive with mass $M>0$. Choose smooth
positive measures $\mu_j$ with the same mass and
$\mu_j\rightharpoonup\mu$. Their mean-zero solutions $q_j$ have
a common lower curvature bound $-M\eta$. Normalized compactness
bounds $\sup_Cq_j$: their mean-zero condition makes this
supremum the negative mean of their supremum-zero normalization.
Hence $(q_j)$ is precompact in $L^1(C)$. Every limit has mean
zero and Hessian $\mu-M\eta$. Two such limits differ by a
harmonic function, hence by a constant, which their means force
to be zero. Therefore $q_j\to q_\mu$ in $L^1(C)$.

By \cref{curve-map:lone-pullback}, $q_j\circ p\to q_\mu\circ p$
in $L^1(X)$. Thus $p^*\mu_j\to p^*\mu$ as currents and their
smooth Hermitian traces converge weakly. For smooth $\mu_j$,
the definition of $S$ gives
$p_*L_\gamma(\mu_j)=\kappa_\gamma\mu_j$. Passing to the limit
proves the identity for $\mu$. The case $M=0$ is zero, and
linearity proves the signed identity.

For a positive $\mu$, split it into $\mathbf1_B\mu$ and
$\mathbf1_{C\setminus B}\mu$. Their traces are positive; the
pushforward identity shows that they are respectively carried
by $p^{-1}(B)$ and its complement. Additivity then gives the
restriction identity. Apply the same argument to the positive
and negative Jordan parts of a signed measure. Their disjoint
Borel carriers pull back to disjoint carriers, so their traces
are the Jordan parts upstairs. The pushforward identity gives
the asserted total-variation norm.

Finally the Poincar\'e--Lelong formula in the book normalization
gives $p^*\delta_s=[F_s]$, including every zero multiplicity of
a local equation of the fiber. Its trace has total mass
$\kappa_\gamma$, by the pushforward identity with $\mu=\delta_s$.
\end{proof}

\subsection{The full model space}

Assume now $n\geq2$ and $\beta>0$, and put $A=\int_C\beta$,
$\theta=p^*\beta$. For $v\in\PSH(C,\beta)$ write
$T_v=\beta+\ddc v$ and $\mu_v=\mathbf1_{\{v=-\infty\}}T_v$.

\begin{lemma}\label{curve-map:model-envelopes}
For $U=v\circ p$,
\[
 P_\theta[U]=P_\beta[v]\circ p,\qquad
 P_\theta[U]_{\mathcal I}=P_\beta[v]_{\mathcal I}\circ p.
\]
Pullback identifies both model spaces and their ideal subspaces.
Every potential in $\PSH(X,\theta)$ has zero extended
Monge--Amp\`ere measure and zero upper and lower singularity volumes.
\end{lemma}
\begin{proof}
By \cref{curve-map:potential-descent,curve-map:ideal-types},
ordinary and ideal competitor families correspond under pullback.
The map $p$ is open: on a line through any point its nonconstant
holomorphic germ is open. Hence upper regularization commutes
with pullback, even at a critical point. These observations
prove both envelope identities and the model assertions.

For the mass assertion, $\max(U,-j)=\max(v,-j)\circ p$. Locally
choose a smooth $a$ with $\ddc a=\beta$. The bounded subharmonic
function $a+\max(v,-j)$ has smooth decreasing local approximants.
Their pullbacks have complex Hessian rank at most one, including
at critical points, so their $n$th powers vanish. Bounded product
continuity proves that every bounded truncation has zero top
product. The non-pluripolar measure is therefore zero. Descent
applies to every representative of every ordinary type, proving
the volume assertion.
\end{proof}

For $U=v\circ p$ define
\[
 \lambda_U^\gamma
  =\mathbf1_{\{U=-\infty\}}(\theta+\ddc U)\wedge\gamma^{n-1},
 \qquad
 \rho_\theta(U,V)=\kappa_\gamma^{-1}
       \|\lambda_U^\gamma-\lambda_V^\gamma\|_{\mathrm{TV}}.
\]

\begin{proposition}\label{curve-map:complete-model-metric}
On the full ordinary model space, $\rho_\theta$ is independent
of $\gamma$ and pullback is an isometry from the complete geodesic
curve model space. Its diameter is $2A$. Its ideal subspace is
closed and geodesically convex, equals the analytic closure,
and has the unique nearest-point, $1$-Lipschitz ideal retraction.
The distance of $v\circ p$ to that subspace is the nonatomic polar
mass of $v$. Metric convergence implies $L^1(X)$ convergence.
The upper-volume comparison is identically zero.
\end{proposition}
\begin{proof}
The Borel restriction identity of \cref{curve-map:gauduchon-trace}
applied to $T_v$ and $\{v=-\infty\}$ gives
$\lambda_U^\gamma=L_\gamma(\mu_v)$. Its variation identity yields
\[
 \rho_\theta(v\circ p,w\circ p)
       =\|\mu_v-\mu_w\|_{\mathrm{TV}}=d_\beta(v,w).
\]
Together with \cref{curve-map:model-envelopes}, this proves
independence and an isometry onto the entire model space.
Completeness and geodesics follow from
\cref{cor:curve-model-completeness}. The mass bound gives diameter
at most $2A$, and the models of $A\delta_s$ and $A\delta_t$ for
distinct base points attain it.

Ideal projection commutes with pullback by the envelope identity.
Thus \cref{prop:curve-analytic-distance-defect} gives its unique
nearest-point property, the Lipschitz estimate and the distance
formula. By \cref{curve-map:analytic-types}, analytic models
correspond in both directions. Every analytic model is ideal by
\cref{prop:analytic-envelope-type}. The analytic curve models of
positive curve mass are dense in the ideal subspace by
\cref{cor:curve-model-completeness}; their pullbacks have zero
top-degree mass. This proves the asserted analytic closure.
Closedness and geodesic convexity follow from those of atomic
polar measures.

Curve metric convergence gives base $L^1$ convergence by
\cref{thm:curve-model-capacity-convergence}, and
\cref{curve-map:lone-pullback} then gives convergence upstairs.
Finally the envelope lemma makes all top masses zero, including
that of each pairwise maximum. This makes the upper-volume
comparison zero.
\end{proof}

\begin{corollary}\label{curve-map:all-obstacles}
For $U=v\circ p$ and $R=P_\theta[U]_{\mathcal I}$, the following
are equivalent: $P_\theta[U]=P_\theta[U]_{\mathcal I}$;
$\mu_v$ is atomic; $U\in\mathcal E_{\mathrm{env}}(X,\theta,R)$;
and the ordinary and ideal envelopes agree below every smooth
real obstacle on $X$.
\end{corollary}
\begin{proof}
Put $r=P_\beta[v]_{\mathcal I}$. The envelope lemma gives
$R=r\circ p$ and makes the first condition equivalent to
$P_\beta[v]=r$. By \cref{cor:curve-arbitrary-background-criterion},
this is equivalent to atomicity and to
$v\in\mathcal E_{\mathrm{env}}(C,\beta,r)$.

For each $b\geq1$, descent and surjectivity give a bijection
between all stretching competitors satisfying
$b^{-1}h+(1-b^{-1})r\leq v$ and their upstairs counterparts.
Only convex combinations occur, so the correspondence remains
valid at poles. It identifies the two relative-class conditions.
Then \cref{prop:ideal-obstacle-reprojection} gives equality below
every smooth obstacle on $X$, including obstacles that do not
descend. Conversely that equality includes the zero obstacle.
\end{proof}
